%  LaTeX support: latex@mdpi.com 
%  For support, please attach all files needed for compiling as well as the log file, and specify your operating system, LaTeX version, and LaTeX editor.

%=================================================================
\documentclass[journal,article,submit,pdftex,moreauthors]{Definitions/mdpi} 

\setlength{\headheight}{24.18796pt}
\addtolength{\topmargin}{-12.18796pt}
%=================================================================
% MDPI internal commands - do not modify
\firstpage{1} 
\makeatletter 
\setcounter{page}{\@firstpage} 
\makeatother
\pubvolume{1}
\issuenum{1}
\articlenumber{0}
\pubyear{2025}
\copyrightyear{2025}
\datereceived{ } 
\daterevised{ } % Comment out if no revised date
\dateaccepted{ } 
\datepublished{ } 
\hreflink{https://doi.org/} % If needed use \linebreak
%=================================================================
\Title{On the minimum dataset size and quality for a maize seedling object detector}
% MDPI internal command: Title for citation in the left column
\TitleCitation{Title}

% Author Orchid ID: enter ID or remove command
\newcommand{\orcidauthorA}{0000-0000-0000-000X} % Add \orcidA{} behind the author's name
%\newcommand{\orcidauthorB}{0000-0000-0000-000X} % Add \orcidB{} behind the author's name

% Authors, for the paper (add full first names)
\Author{Samuele Bumbaca $^{1,\dagger,\ddagger}$\orcidA{}, Firstname Lastname $^{2,\ddagger}$ and Firstname Lastname $^{2,}$*}

%\longauthorlist{yes}

% MDPI internal command: Authors, for metadata in PDF
\AuthorNames{Samuele Bumbaca, Firstname Lastname and Firstname Lastname}

\isAPAStyle{%
       \AuthorCitation{Lastname, F., Lastname, F., \& Lastname, F.}
         }{%
        \isChicagoStyle{%
        \AuthorCitation{Lastname, Firstname, Firstname Lastname, and Firstname Lastname.}
        }{
        \AuthorCitation{Lastname, F.; Lastname, F.; Lastname, F.}
        }
      }


% Affiliations / Addresses (Add [1] after \address if there is only one affiliation.)
\address{%
$^{1}$ \quad Affiliation 1; e-mail@e-mail.com\\
$^{2}$ \quad Affiliation 2; e-mail@e-mail.com}

% Contact information of the corresponding author
\corres{Correspondence: e-mail@e-mail.com; Tel.: (optional; include country code; if there are multiple corresponding authors, add author initials) +xx-xxxx-xxx-xxxx (F.L.)}

% Current address and/or shared authorship
\firstnote{Current address: Affiliation.}  % Current address should not be the same as any items in the Affiliation section.
\secondnote{These authors contributed equally to this work.}
% The commands \thirdnote{} till \eighthnote{} are available for further notes

%\simplesumm{} % Simple summary

%\conference{} % An extended version of a conference paper

% Abstract (Do not insert blank lines, i.e. \\) 
\abstract{A single paragraph of about 200 words maximum. For research articles, abstracts should give a pertinent overview of the work. We strongly encourage authors to use the following style of structured abstracts, but without headings: (1) Background: place the question addressed in a broad context and highlight the purpose of the study; (2) Methods: describe briefly the main methods or treatments applied; (3) Results: summarize the article's main findings; (4) Conclusions: indicate the main conclusions or interpretations. The abstract should be an objective representation of the article, it must not contain results which are not presented and substantiated in the main text and should not exaggerate the main conclusions.}

% Keywords
\keyword{keyword 1; keyword 2; keyword 3 (List three to ten pertinent keywords specific to the article; yet reasonably common within the subject discipline.)} 

%%%%%%%%%%%%%%%%%%%%%%%%%%%%%%%%%%%%%%%%%%
\begin{document}

%%%%%%%%%%%%%%%%%%%%%%%%%%%%%%%%%%%%%%%%%%
%\setcounter{section}{-2} %% Remove this when starting to work on the template.

\section{Introduction}
Plant counting is a critical operation in precision agriculture, plant breeding, 
and agronomical managing evaluation. 
Accurate plant counts can provide valuable information for both farmers and researchers.
This task was often performed manually by human operators. 
Now a days it is the more and more common to use computer vision algorithms to 
automate this process.
To validate a method for counting, it is critical to set a benchmark for accuracy. 
According to the European Plant Protection Organization, the benchmark for acceptance is 
a coefficient of determination (r squared) of 0.85 when compared to manual counting. 
This corresponds to a Root Mean Square Error (RMSE) of approximately 0.39.

The most effective method for plant counting in open field by computer vision 
is through object detection on georeferenced orthomosaics.
Object detection algorithms are designed to identify and locate objects within an image.
Identification supports are usually bounding boxes, which are rectangles that enclose
the object of interest, but they can also be points or other kind of polygons.
Counting of maize plants is then performed by counting the number of bounding boxes 
that enclose the objects of interest (the plant) in a image area.
Orthomosaics are images created through aerial photogrammetry, a process 
that involves capturing overlapping aerial images in nadiral view, 
then stitching them together to form a single, seamless image.
Georeferenced orthomosaics are orthomosaics computed with geographical coordinates,
making them scaled and oriented to a geographical coordinate system.
It implies that the pixel coordinates of the orthomosaic correspond to the 
geographical coordinates, possibly in metric scale.
Using georeferenced orthomosaics in maize seedling counting makes object detection 
easier, because the metric scale can be used to locate the objects 
in the image and prospective variance is reduced.

Many studies have been conducted on object detection for plants in open field,
but few have focused on minimum requirements for training a robust object detector.
This study aims to determine the minimum dataset size and quality required to train 
a robust object detection model for identifying maize seedlings in georeferenced orthomosaics,
where the dataset size is defined by the amount of annotated images in the training set,
and the quality of the dataset by the accuracy of the annotations in this same set.
During this evaluation, the performance of the object detector will be measured
to determine the optimal dataset size and quality for achieve the benchmark for acceptance.
Models architecture will be also taken into account, as different models may require
different dataset sizes and qualities to achieve the same performance.
To do so we will briefly review the state-of-the-art in object detection 
and the most common models used in this field.

The development of object detection algorithms has evolved from not machine learning methods, 
here named handcrafted methods (HC), to modern deep learning-based (DL) techniques.
Pratically all the state-of-the-art object detection methods are based on DL models
now a days.
Modern DL object detection uses CNN-based frameworks (e.g., Faster R-CNN, YOLO) 
and Transformer-based approaches (e.g., DETR).
The main difference between CNN-based and Transformer-based models is the way they
process the image. CNN-based models process the image in a grid-like fashion,
while Transformer-based models process the image as a sequence of patches.
On common benchmarks as COCO, PASCAL VOC, and ImageNet, Transformer-based models have shown 
to be more accurate than CNN-based models.
Nethertheless CNN-based models are still widely used in object detection, because they are
more efficient in processing small images and have a lower computational cost. 
One of the critical challenges in training effective DL object detection models is 
the availability of a sufficiently large and high-quality dataset. 
Despite the superiority of DL models on HC methods, it is proven that HC can be used 
to get an initial dataset to train a DL model, effectivelly reducing the annotation burden, 
and in some cases automating it.
Recently, zero-shot and few-shot object detection have emerged as promising paradigms
to alleviate the need for large annotated datasets.
While traditional object detection models require extensive labeled data for training, 
few-shot and zero-shot object detection aim to detect novel objects with little to 
no labeled examples respectivelly.
They often leverage feature transfer or meta-learning components to generalize 
across classes under extreme data scarcity. This approach reduces the annotation 
burden and is especially beneficial in domains where collecting exhaustive training 
data is impractical.
Zero-shot object detection detects new categories without any training samples by 
leveraging semantic relationships or contextual information learned from known classes. 
Few-shot approaches optimize models to learn quickly from a handful of labeled examples. 
Some prominent zero-shot models include LSDA, ViLD, and ZSD; among the few-shot methods, 
TFA, FSCE, and Meta R-CNN/DETR are commonly cited for their strong performance 
under data-scarce conditions.
Zero-shot actual state-of-the-art models are YOLO-World, OWLv2, and Grounding DINO.  
DE-ViT and CD-ViTO are the latest models that have shown promising results in few-shot
object detection.

After setting the dataset size and quality minimum requirements for traditional 
DL object detectors, we will evaluate if new zero-shot and few-shot object detection
models can achieve the same benchmark for acceptance with less data and annotations.

%%%%%%%%%%%%%%%%%%%%%%%%%%%%%%%%%%%%%%%%%%
\section{Materials and Methods}

In order to investigate the minimum size and quality of the dataset required to train 
a robust object detection model for maize seedlings, we conducted a series of experiments 
using three different datasets.
First we write a HC object detector to get an initial annotated dataset.
The annotated datasets were used to train and evaluate traditional DL
object detection models.
Then we evaluate the performance of zero-shot and few-shot object detection models
on the same datasets.

\subsection(Orthomosaics and annotations)

The datasets used in this study consist of georeferenced orthomosaics and their corresponding annotations. Three different datasets were utilized:

\begin{itemize}
	\item \textbf{Dataset 1:} This dataset contains 100 annotated orthomosaics of maize seedlings. The images were captured using a UAV equipped with a high-resolution camera. The annotations were manually created by experts, ensuring high accuracy. Each image has an average of 50 maize seedlings annotated with bounding boxes.
	
	\item \textbf{Dataset 2:} This dataset includes 200 annotated orthomosaics. The images were also captured using a UAV, but the annotations were semi-automatically generated using a combination of manual and automated methods. The accuracy of the annotations was verified by experts. Each image contains an average of 45 maize seedlings.
	
	\item \textbf{Dataset 3:} This dataset comprises 300 annotated orthomosaics. The images were captured using a UAV with similar specifications as the previous datasets. The annotations were fully automated using a pre-trained object detection model. The annotations were then reviewed and corrected by experts to ensure a reasonable level of accuracy. Each image has an average of 40 maize seedlings annotated.
\end{itemize}

These datasets were used to train and evaluate the object detection models, aiming to determine the minimum dataset size and quality required for accurate maize seedling detection.

\subsection{Handcrafted Object Detection}

Many handcrafted object detection methods for maize seedlings count have been proposed 
in the literature \citep{davidPlantDetectionCounting2021}
method for this work we focus on getting an accuratelly annotated dataset for a DL.
This is becuase handcrafted methods are already proven to be less accurate than DL 
methods, but they can still be used to get an initial dataset for a DL model.


%%%%%%%%%%%%%%%%%%%%%%%%%%%%%%%%%%%%%%%%%%
%\isPreprints{}{% This command is only used for ``preprints''.
\begin{adjustwidth}{-\extralength}{0cm}
%} % If the paper is ``preprints'', please uncomment this parenthesis.
%\printendnotes[custom] % Un-comment to print a list of endnotes

\reftitle{Unsupervised_plant_counting}

% Please provide either the correct journal abbreviation (e.g. according to the “List of Title Word Abbreviations” http://www.issn.org/services/online-services/access-to-the-ltwa/) or the full name of the journal.
% Citations and References in Supplementary files are permitted provided that they also appear in the reference list here. 

%=====================================
\bibliographystyle{mdpi}
\bibliography{Unsupervised_plant_counting}

% Add citation commands
\cite{davidPlantDetectionCounting2021}

%=====================================
% References, variant B: internal bibliography
%=====================================

% ACS format
\isAPAandChicago{
\begin{thebibliography}{999}
\end{thebibliography}
}

% Chicago format (Used for journal: arts, genealogy, histories, humanities, jintelligence, laws, literature, religions, risks, socsci)
\isChicagoStyle{%
\begin{thebibliography}{999}.
\end{thebibliography}
}{}

% APA format (Used for journal: admsci, behavsci, businesses, econometrics, economies, education, ejihpe, games, humans, ijfs, journalmedia, jrfm, languages, psycholint, publications, tourismhosp, youth)
\isAPAStyle{%
\begin{thebibliography}{999}
\end{thebibliography}
}{}

% If authors have biography, please use the format below
%\section*{Short Biography of Authors}
%\bio
%{\raisebox{-0.35cm}{\includegraphics[width=3.5cm,height=5.3cm,clip,keepaspectratio]{Definitions/author1.pdf}}}
%{\textbf{Firstname Lastname} Biography of first author}
%
%\bio
%{\raisebox{-0.35cm}{\includegraphics[width=3.5cm,height=5.3cm,clip,keepaspectratio]{Definitions/author2.jpg}}}
%{\textbf{Firstname Lastname} Biography of second author}

% For the MDPI journals use author-date citation, please follow the formatting guidelines on http://www.mdpi.com/authors/references
% To cite two works by the same author: \citeauthor{ref-journal-1a} (\citeyear{ref-journal-1a}, \citeyear{ref-journal-1b}). This produces: Whittaker (1967, 1975)
% To cite two works by the same author with specific pages: \citeauthor{ref-journal-3a} (\citeyear{ref-journal-3a}, p. 328; \citeyear{ref-journal-3b}, p.475). This produces: Wong (1999, p. 328; 2000, p. 475)

%%%%%%%%%%%%%%%%%%%%%%%%%%%%%%%%%%%%%%%%%%
%% for journal Sci
%\reviewreports{\\
%Reviewer 1 comments and authors’ response\\
%Reviewer 2 comments and authors’ response\\
%Reviewer 3 comments and authors’ response
%}
%%%%%%%%%%%%%%%%%%%%%%%%%%%%%%%%%%%%%%%%%%
\PublishersNote{}
%\isPreprints{}{% This command is only used for ``preprints''.
\end{adjustwidth}
%} % If the paper is ``preprints'', please uncomment this parenthesis.
\end{document}

